% Part: methods
% Chapter: proofs
% Section: proof-by-contradiction

\documentclass[../../../include/open-logic-section]{subfiles}

\begin{document}

\olfileid{mth}{prf}{con}

\olsection{ٹکراؤ توں ثبوت}

پہلاں ایہ ویکھو کہ ٹکراؤ توں ثبوت، استنباط دا اوہ ڈھنگ اے جیہڑا نفی والے دعوے ثابت کرن لئی ورتیا جاندا اے۔ منّ لو تسی وکھاؤنا چاہندے او کہ کوئی دعویٰ~$p$ \emph{غلط} اے، یعنی تسی~$\lnot p$ وکھاؤنا چاہندے او۔ سبھ توں آس والا طریقہ ایہ اے کہ (ا) $p$ نوں سچ منّو، تے (ب) وکھاؤ کہ ایس منّی ہوئی گل توں کوئی ایہو جیہی گل نکلدی اے جیہڑی تہانوں پتا اے کہ غلط اے۔ «ایہو جیہی گل جیہڑی پتا اے کہ غلط اے» کوئی ایہو جیہا نتیجہ ہو سکدا اے جیہڑا $p$ آپ نال، یا سارے دعوے دی کسے ہور مفروضہ شرط نال ٹکراؤندا اے۔ مثلاً «جے $q$ تاں $\lnot p$» دے ثبوت وچ $q$ نوں سچ منّنا تے اوہتوں~$\lnot
p$ ثابت کرنا پیندا اے۔ جے تسی $\lnot p$ ٹکراؤ توں ثابت کردے او، تاں $p$ وی منّنا پئے گا،~$q$ دے نال۔ جے تسی $\lnot q$، $p$ توں ثابت کر سکو، تاں تسی وکھا دتا کہ مفروضہ~$p$ توں ایہو جیہی گل نکلدی اے جیہڑی تہاڈے دوجے مفروضے~$q$ نال ٹکراؤندی اے، کیوں جے $q$~تے $\lnot q$ دوویں سچ نہیں ہو سکدے۔ بے شک ٹکراؤ ثابت کردیاں تہانوں استنباط دے ہور ڈھنگ ورتنے پَین گے تے تعریفاں وی کھولھنیاں پَین گئیاں۔ اک مثال ویکھیے۔

\begin{prop}
  جے $A \subseteq B$ تے $B = \emptyset$، تاں $A$ دے !!{element}s کوئی نہیں۔
\end{prop}

\begin{proof}
  منّ لو $A \subseteq B$ تے $B = \emptyset$۔ اسیں وکھاؤنا چاہندے آں کہ $A$ دے !!{element}s کوئی نہیں۔
  \begin{quote}
    ایہ شرطیہ دعویٰ اے، سو مقدم منّ کے تالی ثابت کرنا چاہندے آں۔ تالی ایہ اے کہ $A$ دے !!{element}s کوئی نہیں۔ ایہ گل ذرا کھولھ کے آکھ سکدے آں کہ کوئی ایہو جیہا~$x \in A$ نہیں۔
  \end{quote}
  $A$ دے !!{element}s کوئی نہیں اُدوں تے صرف اُدوں جدوں کوئی ایہو جیہا~$x$ نہیں جیہدے لئی $x \in A$۔
  \begin{quote}
    سو اسیں طے کیتا اے کہ ثابت کرن والی گل اصل وچ نفی والا دعویٰ~$\lnot p$ اے، یعنی کوئی $x \in A$ نہیں۔ ٹکراؤ توں ثبوت لئی سانوں اوہدے نال ملدا اثباتی دعویٰ~$p$ منّنا پئے گا، یعنی کوئی $x \in A$ اے، تے اوہتوں ٹکراؤ ثابت کرنا پئے گا۔ ایہ دسّن لئی کہ اسیں ٹکراؤ توں ثبوت کردے آں، «ٹکراؤ کڈھن لئی منّ لو» یا بس «اُلٹ منّ لو» لکھدے آں تے فیر مفروضہ~$p$ دس دیندے آں۔
  \end{quote}
  اُلٹ منّ لو، کوئی $x \in A$ اے۔
  \begin{quote}
    ہُن ایہ اوہ نواں مفروضہ اے جیہڑا اسیں ٹکراؤ کڈھن لئی ورتاں گے۔ ساڈے دو ہور مفروضے نیں، $A \subseteq B$ تے $B = \emptyset$۔ پہلے توں ملدا اے کہ $x \in B$:
  \end{quote}
  کیوں جے $A \subseteq B$، سو $x \in B$۔
  \begin{quote}
    پر کیوں جے $B = \emptyset$، سو $B$ دا ہر !!{element} (جیویں $x$) لازماً~$\emptyset$ دا وی !!a{element} اے۔
  \end{quote}
  کیوں جے $B = \emptyset$، سو $x \in \emptyset$۔ ایہ ٹکراؤ اے، کیوں جے تعریف مطابق $\emptyset$ دے !!{element}s کوئی نہیں۔
  \begin{quote}
    ایس نال ثبوت پہلاں ای پورا اے۔ اسیں اپنے مفروضیاں توں لوڑیندی گل (اک ٹکراؤ) تک پہنچ گئے آں، تے ایس دا مطلب اے کہ سارے مفروضے سچ نہیں ہو سکدے۔ پہلے دو مفروضیاں ($A \subseteq B$ تے $B = \emptyset$) اُتے کوئی اعتراض نہیں۔ سو آخر وچ پیش کیتا مفروضہ (کوئی $x \in A$ اے) غلط ہونا پئے گا۔ پر جے اسیں ثبوت نوں ہور صاف رکھنا چاہیئے، تاں ایہ گل کھولھ کے آکھ سکدے آں۔
  \end{quote}
  سو ساڈا مفروضہ کہ کوئی $x \in A$ اے لازماً غلط اے۔ ایس لئی ٹکراؤ توں ثبوت مطابق $A$ دے !!{element}s کوئی نہیں۔
\end{proof}

ہر اثباتی دعویٰ سِدھے ای اک نفی والے دعوے دے برابر اے۔ $p$ اُدوں تے صرف اُدوں جدوں $\lnot\lnot p$۔ سو ٹکراؤ توں ثبوت اثباتی دعوے «بالواسطہ» ثابت کرن لئی وی ورتیا جا سکدا اے۔~$p$ ثابت کرن لئی اینوں نفی والا دعویٰ $\lnot\lnot p$ سمجھو۔ جے اسیں $\lnot p$ توں ٹکراؤ ثابت کر سکیے، تاں ٹکراؤ توں ثبوت نال $\lnot\lnot p$ ثابت ہو گیا، تے ایس لئی~$p$ وی۔

پچھلی مثال وچ اسیں نفی والا دعویٰ ثابت کرنا چاہندے ساں، یعنی $A$ دے !!{element}s کوئی نہیں۔ سو ٹکراؤ توں ثبوت لئی جیہڑا مفروضہ منّیا سی (یعنی کوئی $x \in A$ اے) اوہ اثباتی دعویٰ سی۔ ایس نال سانوں دلیل لئی اک گل ملی، یعنی منّیا ہویا $x \in A$، جیہدے بارے دلیل دے کے اسیں $x \in \emptyset$ تک پہنچے۔

اثباتی دعوے دا بالواسطہ ثبوت کردیاں ٹکراؤ کڈھن لئی تہاڈا مفروضہ نفی والا ہووے گا۔ پر اکثر تسی اثباتی دعوے نوں سوکھ نال نفی والی صورت وچ، تے نفی والے دعوے نوں اثباتی صورت وچ فیر لکھ سکدے او۔ جے اسیں «$A = \emptyset$» ثابت کردے، نفی والی تالی «$A$ دے !!{element}s کوئی نہیں» دی تھاں، تاں ساڈا پچھلا ثبوت بنیادی طور اُتے اوہو رہنا سی۔ ($=$ دی تعریف مطابق «$A = \emptyset$» اک عام دعویٰ اے، کیوں جے کھولھ کے ایہدا مطلب اے «ہر !!{element} جیہڑا~$A$ دا اے، اوہ~$\emptyset$ دا وی !!{element} اے، تے اُلٹے پاسے وی»۔) پر ایہ سوکھ نال ویکھ سکدے آں کہ ایہ نفی والے دعوے «ایہ گل نہیں کہ کوئی $x \in A$ اے» دے برابر اے۔

سو کدی $\lnot p$ نوں مفروضہ منّ کے کم کرنا،~$p$ نوں سِدھا ثابت کرن نالوں سوکھا اے۔ بھاویں سِدھا ثبوت اوہنا ای سوکھا ہووے یا ہور وی سوکھا (جیویں اگلیاں مثالاں وچ)، کجھ لوک بالواسطہ طریقہ پسند کردے نیں۔ جے دوہری نفی نال تہانوں بھلیکھا پئے، تاں کسے دعوے دے ٹکراؤ توں ثبوت نوں اوہدے \emph{اُلٹ} دعوے توں ٹکراؤ دا ثبوت سمجھو۔ سو $\lnot
p$ دا ٹکراؤ توں ثبوت مفروضے~$p$ توں ٹکراؤ دا ثبوت اے، تے~$p$ دا ٹکراؤ توں ثبوت~$\lnot p$ توں ٹکراؤ دا ثبوت اے۔

\begin{prop}
$A \subseteq A \cup B$.
\end{prop}

\begin{proof}
  اسیں وکھاؤنا چاہندے آں کہ $A \subseteq A \cup B$۔
  \begin{quote}
    پہلاں ویکھن نوں ایہ اثباتی دعویٰ اے۔ ہر $x \in A$، $A \cup B$ وچ وی اے۔ اوہدی نفی ایہ اے کہ کوئی $x \in A$، $\notin A \cup B$ اے۔ سو اسیں ایس نفی والے دعوے نوں منّ کے تے اوہتوں ٹکراؤ وکھا کے دعویٰ بالواسطہ ثابت کر سکدے آں۔
    \end{quote}
  اُلٹ منّ لو، یعنی $A \nsubseteq A \cup B$۔
  \begin{quote}
    ساڈے کول $A \subseteq A \cup B$ دی تعریف اے، ہر $x \in A$، $\in A \cup B$ وی اے۔ $A \nsubseteq A \cup B$ دا مطلب سمجھن لئی سانوں کھولھی ہوئی تعریف اُتے بنیادی منطقی کم کرنا پئے گا۔ ایہ غلط اے کہ ہر $x \in A$، $\in
    A \cup B$ وی اے، اُدوں تے صرف اُدوں جدوں \emph{کوئی}~$x \in A$ اے جیہڑا $\notin C$ اے۔ (ایتھے بڑا دھیان رکھو۔ کئی طالب علماں دے ٹکراؤ توں ثبوت دی کوشش ایس لئی نہیں چل دی کہ اوہ «سارے $A$، $B$ نیں» ورگے دعوے دی نفی غلط سمجھدے نیں۔) دوجے لفظاں وچ، $A \nsubseteq A \cup B$ اُدوں تے صرف اُدوں جدوں کوئی $x$ اے جیہدے لئی $x \in A$ تے $x \notin A \cup B$۔ ایتھوں اگے گل سوکھی اے۔
  \end{quote}
  سو کوئی $x \in A$ اے جیہدے لئی $x \notin A \cup B$۔ $\cup$ دی تعریف مطابق $x \in A \cup B$ اُدوں تے صرف اُدوں جدوں $x \in A$ یا $x \in
  B$۔ کیوں جے $x \in A$، سو $x \in A \cup B$۔ ایہ مفروضے $x \notin A \cup B$ نال ٹکراؤندا اے۔
\end{proof}

\begin{prob}
\emph{بالواسطہ} ثابت کرو کہ $A \cap B \subseteq A$۔
\end{prob}

\begin{prop}
جے $A \subseteq B$ تے $B \subseteq C$، تاں $A \subseteq C$۔
\end{prop}

\begin{proof}
  منّ لو $A \subseteq B$ تے $B \subseteq C$۔ اسیں وکھاؤنا چاہندے آں کہ $A
  \subseteq C$۔
  \begin{quote}
    آؤ بالواسطہ طریقہ ورتیے۔ جیہڑی گل ثابت کرنا چاہندے آں، اوہدی نفی منّ لیندے آں۔
  \end{quote}
  اُلٹ منّ لو، یعنی $A \nsubseteq C$۔
  \begin{quote}
    پہلاں وانگ، اسیں دلیل دیندے آں کہ $A \nsubseteq C$ اُدوں تے صرف اُدوں جدوں ایہ گل نہیں کہ ہر $x \in A$، $\in C$ وی اے، یعنی کوئی $x \in A$، $\notin C$ اے۔ فکر نہ کرو، مشق نال ایہو جیہیاں نفیاں کھولھن لئی بوہتا سوچنا نہیں پئے گا۔
  \end{quote}
  دوجے لفظاں وچ، کوئی~$x$ اے جیہدے لئی $x \in A$ تے $x \notin C$۔
  \begin{quote}
    ہُن ایہ گل ٹکراؤ تک پہنچن لئی ورت سکدے آں۔ بے شک ایس کم لئی دوویں ہور مفروضے وی ورتنے پَین گے۔
  \end{quote}
  کیوں جے $A \subseteq B$، سو $x \in B$۔ کیوں جے $B \subseteq C$، سو $x \in
  C$۔ پر ایہ $x \notin C$ نال ٹکراؤندا اے۔
\end{proof}

\begin{prop}
جے $A \cup B = A \cap B$، تاں $A = B$۔
\end{prop}

\begin{proof}
  منّ لو $A \cup B = A \cap B$۔ اسیں وکھاؤنا چاہندے آں کہ $A = B$۔
  \begin{quote}
    ہُن شروع دی گل جانی پچھانی اے:
  \end{quote}
  ٹکراؤ کڈھن لئی منّ لو کہ $A \neq B$۔
  \begin{quote}
    ٹکراؤ توں ثبوت لئی ساڈا مفروضہ $A \neq
    B$ اے۔ کیوں جے $A = B$ اُدوں تے صرف اُدوں جدوں $A \subseteq B$ تے $B \subseteq A$، ایس توں ملدا اے کہ $A \neq B$ اُدوں تے صرف اُدوں جدوں $A \nsubseteq B$ \emph{یا} $B \nsubseteq
    A$۔ (ویکھو نفیاں نال کم کردیاں دھیان رکھنا کِنّا ضروری اے!{}) ایس فصل توں ٹکراؤ ثابت کرن لئی صورتاں وکھ کر کے ثبوت دیندے آں تے وکھاؤندے آں کہ ہر صورت وچ ٹکراؤ نکلدا اے۔
  \end{quote}
  $A \neq B$ اُدوں تے صرف اُدوں جدوں $A \nsubseteq B$ یا $B \nsubseteq A$۔ اسیں صورتاں وکھ کردے آں۔
  \begin{quote}
    پہلی صورت وچ اسیں $A \nsubseteq B$ منّدے آں، یعنی کسے $x$ لئی، $x \in A$ پر $\notin B$۔ $A \cap B$ دی تعریف اوہناں !!{element}s نال دتی اے جیہڑے $A$ تے $B$ دے سانجھے نیں۔ سو جے کوئی شے کسے اک وچ نہیں، تاں سانجھے حصے وچ نہیں۔ $A \cup B$، $A$ تے $B$ دا اکٹھ اے، سو جیہڑی شے کسے اک وچ اے، اوہ اتحاد وچ وی اے۔ ایس توں پتا لگدا اے کہ $x \in A \cup B$ پر $x \notin A \cap
    B$، تے ایس لئی $A \cap B \neq A \cup B$۔
  \end{quote}

  صورت ۱۔ $A \nsubseteq B$۔ تاں کسے $x$ لئی، $x \in A$ پر $x \notin
  B$۔ کیوں جے $x \notin B$، سو $x \notin A \cap B$۔ کیوں جے $x \in A$، سو $x \in A \cup B$۔ ایس لئی $A \cap B \neq A \cup B$، جیہڑا مفروضے $A \cap B = A \cup B$ نال ٹکراؤندا اے۔

  صورت ۲۔ $B \nsubseteq A$۔ تاں کسے $y$ لئی، $y \in B$ پر $y \notin
  A$۔ پہلاں وانگ، سانوں $y \in A \cup B$ پر $y \notin A \cap B$ ملدا اے، سو $A \cap B \neq A \cup B$، جیہڑا فیر $A \cap B = A \cup
  B$ نال ٹکراؤندا اے۔
\end{proof}

\end{document}
